\documentclass[10pt,a4paper]{amsart}
\usepackage[margin=19mm]{geometry}
\usepackage{amsmath,amssymb,mathtools}
\usepackage[scr=rsfs,cal=euler]{mathalfa}
\usepackage{xfrac}
\usepackage[unicode,hidelinks,bookmarks=false]{hyperref}
\newcommand{\ie}{i.e.\ }
\newcommand{\cf}{cf.\ }
\newcommand{\resp}{respectively\ }
\newcommand{\Subsection}{\subsection}
\providecommand{\nobreakdash}{\nobreak\hbox{-}}
\setlength{\emergencystretch}{2em}
\multlinegap0pt
\usepackage{amssymb}
\usepackage{xcolor}
\usepackage{tikz}
\newtheorem{lemma}{Lemma}
\title{Neural collapse in the orthoplex regime}
\author{James Alcala, Rayna Andreeva, Vladimir A. Kobzar, Dustin G. Mixon, Sanghoon Na, Shashank Sule, Yangxinyu Xie}
\date{Source: arXiv:2603.20587v1 (CC BY 4.0)}
\begin{document}
\maketitle

\noindent\emph{Source and licence.} Neural collapse in the orthoplex regime. Version arXiv:2603.20587v1 (CC BY 4.0), distributed by arXiv under the Creative Commons Attribution 4.0 International licence, http://creativecommons.org/licenses/by/4.0/, as recorded in the arXiv licence metadata for this exact version and shown on the abstract page https://arxiv.org/abs/2603.20587v1. Full licence notice: \texttt{licenses/arxiv-2603.20587-CC-BY-4.0.txt}.\par\smallskip

\noindent\textit{Editorial scope.} Counted unit: the article's lemma lem.concave-convex (source0.tex lines 888--898). The article's complete original proof of it is delivered with it (source0.tex lines 984--1075, one \texttt{proof} environment of 92 lines, which the article places away from the statement). No mathematics is written, repaired or re-derived: the statement and the proof are the article's own lines, in the article's own notation, and no retained line is substituted or re-flowed.  No further source-local context is needed: every cross-reference of every retained line resolves inside the two delivered spans.  Nothing was omitted from any retained span, so the package carries no omission record.  No retained line contains a citation, so the wrapper carries no bibliography and no bibliography entry.  No retained line contains a figure environment, an image-inclusion command or drawing code, so no image is delivered and the package declares no asset dependency.  Editorial formatting only, in addition: an AMS \texttt{amsart} preamble that restates the article's own declarations and macros used by the retained lines, namely the environments \texttt{lemma} on separate counters, together with the standard packages those lines need.  The retained environments print in this document as Lemma 1 in order of appearance, whereas the article prints the same items with the numbering of its own full document.  The complete native source of the article as distributed in the arXiv e-print for this exact version is bundled unchanged as \texttt{sources/196861/source0.tex}.
\begin{lemma}
\label{lem.concave-convex}
For every integer $n>1$, there exist thresholds 
\[
0
<\tau_-(n)
<\tau_+(n)
<\infty
\]
such that for every $\tau<\tau_-(n)$ (resp.\ $\tau>\tau_+(n)$), it holds that $f_{n,\tau}$ is strictly concave (resp.\ convex).
\end{lemma}

\begin{proof}[Proof of Lemma~\ref{lem.concave-convex}]
For convenience, we put
\[
g(x)
:=n-x-1+e^{\beta}+xe^{-\beta/x},
\]
where $\beta:=1/\tau$.
Notably, $g(x)>0$ for every $x\in[1,n-1]$.
We have
\[
f''_{n,\tau}(x)
=\frac{1}{g(x)}\bigg((x+1)g''(x)+2g'(x)-(x+1)\frac{g'(x)^2}{g(x)}\bigg).
\]
It is sometimes convenient to write $u:=1/(\tau x)=\beta/x$.
We have
\[
g'(x)
=-1+(1+u)e^{-u},
\qquad
g''(x)
=\frac{u^2e^{-u}}{x}.
\]
In what follows, we apply the following primitives without further explanation:
\[
1\leq x\leq n-1\,,
\qquad
0\leq u\leq \beta\,.
\]

For the concavity claim, it suffices to show
\[
Q
:=(x+1)g''(x)+2g'(x)-(x+1)\frac{g'(x)^2}{g(x)}
<0.
\]
For the first two terms, we have
\begin{align*}
(x+1)g''(x)
&=(1+\tfrac{1}{x})u^2e^{-u}
\leq 2\beta^2e^{-\beta/n},\\[0.5em]
2g'(x)
&=2\big(-1+(1+u)e^{-u}\big)
\leq2\big(-1+(1+\beta)e^{-\beta/n}\big),
\end{align*}
and since
\[
g(x)
=n-x-1+e^{\beta}+xe^{-\beta/x}
\geq e^{\beta}+xe^{-\beta/x}
> 0,
\]
the third term satisfies
\[
-(x+1)\frac{g'(x)^2}{g(x)}
\leq0.
\]
Putting everything together,
\[
Q
\leq 2\big(-1+(1+\beta+\beta^2)e^{-\beta/n}\big),
\]
which is negative for all large $\beta$, i.e., it suffices to take $\tau<\tau_-(n)$ for some sufficiently small $\tau_-(n)$.

For the convexity claim, it suffices to show $Q>0$.
For the first two terms, we have
\[
(x+1)g''(x)
=(1+\tfrac{1}{x})u^2e^{-u}
\geq(1+\tfrac{1}{n})u^2e^{-u},
\qquad
2g'(x)
=2\big(-1+(1+u)e^{-u}\big).
\]
For the third term, we have
\[
g'(x)^2
=\big(-1+(1+u)e^{-u}\big)^2,
\qquad
\frac{g(x)}{x+1}
=\frac{n+e^\beta}{x+1}-1+\frac{x}{x+1}e^{-u}
\geq\frac{1}{n}+\frac{1}{2}e^{-u}.
\]
Putting everything together,
\begin{align*}
Q
&\geq(1+\tfrac{1}{n})u^2e^{-u}+2\big(-1+(1+u)e^{-u}\big)-\frac{\big(-1+(1+u)e^{-u}\big)^2}{\frac{1}{n}+\frac{1}{2}e^{-u}}\\
&=\tfrac{1}{n}u^2+(\tfrac{1}{3}+\tfrac{1}{n})u^3+\cdots,
\end{align*}
where the last step is the Taylor series expansion with respect to $u$, centered at $0$.
By an application of Taylor's theorem, it follows that $Q>0$ for all sufficiently small $u$.
Since $u=1/(\tau x)\leq 1/\tau$, it suffices to take $\tau>\tau_+(n)$ for some sufficiently large $\tau_+(n)$.
\end{proof}

\end{document}
